\thispagestyle{empty}
\newgeometry{margin=1cm}
\footnotesize

\begin{landscape}
\begin{center}
\section*{Cheatsheet}
\end{center}
\begin{multicols*}{4}

\begin{tcolorboxTemplate}[title=Kolonnen auswählen: \texorpdfstring{\lstinline|SELECT|}{SELECT},colbacktitle=white, coltitle=black]
Auswahl einer oder mehrerer Kolonnen:
\begin{lstlisting}
SELECT *
FROM users
\end{lstlisting}
\end{tcolorboxTemplate}

\begin{tcolorboxTemplate}[title=Dublikate entfernen: \texorpdfstring{\lstinline|SELECT DISTINCT|}{SELECT DISTINCT},colbacktitle=white, coltitle=black]
\begin{lstlisting}
SELECT DISTINCT gender
FROM users
\end{lstlisting}
\end{tcolorboxTemplate}

\begin{tcolorboxTemplate}[title=Zeilen auswählen: \texorpdfstring{\lstinline|WHERE|}{WHERE},colbacktitle=white, coltitle=black]
Beschränkung auf ausgewählte Zeilen:
\begin{lstlisting}
SELECT name, city, gender
FROM users
WHERE gender="male"
\end{lstlisting}

\begin{table}[H]
\footnotesize
  \centering
  \begin{tabular}{c}
    \textbf{Operatoren} \\\midrule
    \texttt{=}                                    \\
    \texttt{<>}                                   \\
    \texttt{<}                                    \\
    \texttt{<=}                                   \\
    \texttt{>}                                    \\
    \texttt{>=}                                   \\
    \lstinline|BETWEEN x AND y|        \\
    \lstinline|IN (wert1, wert2, ...)| \\
    \lstinline|IS NULL|                \\
  \end{tabular}
\end{table}
\end{tcolorboxTemplate}

\begin{tcolorboxTemplate}[title=Ungefähre Treffer: \texorpdfstring{\lstinline|LIKE|}{LIKE},colbacktitle=white, coltitle=black]
Mit folgendem Befehl erhalten wir alle Städte, die mit ``Be...'' beginnen:
\begin{lstlisting}
SELECT DISTINCT city 
FROM users 
WHERE city LIKE "Be%"
\end{lstlisting}
\end{tcolorboxTemplate}

\begin{tcolorboxTemplate}[title=Sortieren: \texorpdfstring{\lstinline|ORDER BY|}{ORDER BY},colbacktitle=white, coltitle=black]
Mit folgendem Befehl erhalten wir alle Städte alphabetisch sortiert:
\begin{lstlisting}
SELECT DISTINCT city 
FROM users 
ORDER BY city DESC
\end{lstlisting}
\begin{itemize}
\item \lstinline|DESC| = absteigend (\textit{descending})
\item \lstinline|ASC| = aufsteigend (\textit{ascending})
\end{itemize}
\end{tcolorboxTemplate}

\begin{tcolorboxTemplate}[title=Aggregatsfunktionen,colbacktitle=white, coltitle=black]
Mögliche Aggregatsfunktionen: \lstinline|COUNT, MAX, MIN, SUM, AVG LENGTH|
\end{tcolorboxTemplate}

\begin{tcolorboxTemplate}[title=Gruppieren: \texorpdfstring{\lstinline|GROUP BY|}{GROUP BY},colbacktitle=white, coltitle=black]
Anzahl Mitglieder pro Stadt:
\begin{lstlisting}
SELECT city, COUNT(*) AS "Mitglieder pro Stadt" 
FROM users 
GROUP BY city
\end{lstlisting}
\end{tcolorboxTemplate}

\begin{tcolorboxTemplate}[title=Filtern nach Gruppieren: \texorpdfstring{\lstinline|HAVING|}{HAVING},colbacktitle=white, coltitle=black]
Durchschnittliche Körpergrösse aller Mitglieder in jeder Stadt berechnen, danach mit \lstinline|HAVING| Restultate beschränken auf Städte, in denen die Menschen durchschnittlich zwischen 150 und 155 gross sind.
\begin{lstlisting}
SELECT city, AVG(centimeters) AS "Durchschnittliche Körpergrösse"
FROM users
GROUP BY city
HAVING `Durchschnittliche Körpergrösse` BETWEEN 150 AND 155
\end{lstlisting}
\end{tcolorboxTemplate}

\begin{tcolorboxTemplate}[title=Erste / Letzte Zeilen: \texorpdfstring{\lstinline|LIMIT|}{LIMIT},colbacktitle=white, coltitle=black]
Drei Städte mit den meisten Mitgliedern:
\begin{lstlisting}
SELECT city, COUNT(*) AS "Anzahl Mitglieder"
FROM users
GROUP BY city
ORDER BY `Anzahl Mitglieder` DESC
LIMIT 3
\end{lstlisting}
\end{tcolorboxTemplate}

\begin{tcolorboxTemplate}[title=Verzweigungen: \texorpdfstring{\lstinline|CASE WHEN|}{CASE WHEN},colbacktitle=white, coltitle=black]
Für alle Männer berechnen, ob sie grösser, kleiner oder gleich dem Durchschnitt sind:

\begin{lstlisting}
SELECT name, centimeters,
CASE 
	WHEN centimeters > 179 THEN 'Gross'
    WHEN centimeters < 179 THEN 'Klein'
    ELSE 'Genau im Durchschnitt'
END AS Durchschnittlich
FROM users
WHERE gender="male"
\end{lstlisting}
\end{tcolorboxTemplate}

\begin{tcolorboxTemplate}[title=Mehrere Tabellen verbinden: \texorpdfstring{\lstinline|JOIN|}{JOIN},colbacktitle=white, coltitle=black]
Falls \lstinline|Kolonne_ID| in beiden Tabellen vorhanden:
\begin{lstlisting}
SELECT kolonne1, kolonne2, ...
FROM tabelle1 JOIN tabelle2 USING (Kolonne_ID)
\end{lstlisting}

Ansonsten:
\begin{lstlisting}
SELECT kolonne1, kolonne2, ...
FROM tabelle1 JOIN tabelle2 
ON tabelle1.kolonne_id1 = tabelle2.kolonne_id2
\end{lstlisting}
\end{tcolorboxTemplate}

\begin{tcolorboxTemplate}[title=Tabellen erstellen: \texorpdfstring{\lstinline|CREATE table|}{CREATE table},colbacktitle=white, coltitle=black]
\begin{lstlisting}
CREATE TABLE tabellenname(
  kolonne1 eigenschaften1,
  kolonne2 eigenschaften2,
  ...
)
\end{lstlisting}
\end{tcolorboxTemplate}

\begin{tcolorboxTemplate}[title=Tabellen löschen: \texorpdfstring{\lstinline|DROP table|}{DROP table},colbacktitle=white, coltitle=black]
\begin{lstlisting}
DROP TABLE tabellenname
\end{lstlisting}

\end{tcolorboxTemplate}

\begin{tcolorboxTemplate}[title=Daten einfügen: \texorpdfstring{\lstinline|INSERT|}{INSERT},colbacktitle=white, coltitle=black]
\begin{lstlisting}
INSERT INTO tabelle_name (kolonne1, kolonne2, kolonne3, ...)
VALUES (wert1, wert2, wert3, ...)
\end{lstlisting}
\end{tcolorboxTemplate}

\begin{tcolorboxTemplate}[title=Daten verändern: \texorpdfstring{\lstinline|UPDATE|}{UPDATE},colbacktitle=white, coltitle=black]
\begin{lstlisting}
UPDATE tabelle_name
SET kolonne1 = wert1, kolonne2=wert2, ...
WHERE bedingung(en)
\end{lstlisting}
\end{tcolorboxTemplate}

\begin{tcolorboxTemplate}[title=Einträge löschen: \texorpdfstring{\lstinline|DELETE FROM|}{DELETE FROM},colbacktitle=white, coltitle=black]
\begin{lstlisting}
DELETE FROM users
WHERE username="guenther37"
\end{lstlisting}
\end{tcolorboxTemplate}

\end{multicols*}
\end{landscape}
